\documentclass[11pt]{article}

\usepackage[T1]{fontenc}
\usepackage{amsmath,amssymb,amsthm}
\usepackage{hyperref}

\title{Misspecification-Robust Inference for the Incremental Fit Indices\\
CFI and TLI under an Estimated Weight}
\author{}
\date{\today}

\newcommand{\E}{\mathbb{E}}
\newcommand{\tr}{\operatorname{tr}}
\newcommand{\Var}{\operatorname{Var}}
\newcommand{\Cov}{\operatorname{Cov}}
\newcommand{\vech}{\operatorname{vech}}
\newcommand{\cfi}{\widehat{\mathrm{CFI}}}
\newcommand{\tli}{\widehat{\mathrm{TLI}}}
\newcommand{\dto}{\xrightarrow{d}}
\newcommand{\pto}{\xrightarrow{p}}

\theoremstyle{plain}
\newtheorem{theorem}{Theorem}
\newtheorem{lemma}{Lemma}
\newtheorem{corollary}{Corollary}
\theoremstyle{remark}
\newtheorem*{remark}{Remark}

\begin{document}
\maketitle

\section*{Purpose}

The companion notes fix the estimated-weight inference of a \emph{single}
quadratic-form criterion. \texttt{rmsea\_asymptotics.tex} and the RMSEA path of
\texttt{higher\_order\_discrepancy\_misspec.tex} treat the discrepancy
$F=r'Wr$ itself; the CRMR path treats the correlation residual sum $G=r'V_0r$.
Both reduce to one object: a statistic $T=N\hat F$, a bias-corrected
noncentrality $\hat\delta=T-\tr(Q\Gamma_x)$ where $Q$ is a profile Hessian and
$\Gamma_x$ the joint NACOV of the extended first-stage vector
$x=(u,\gamma)$, and a sampling variance $\Var(T)=N\,g'\Gamma_x g$ obtained by
the delta method from an envelope-score gradient $g$.

This note opens the program's \emph{first incremental} (two-model) index. CFI
and TLI are not functions of one criterion; they compare the fitted user model
to the independence baseline,
\begin{equation}\label{eq:cfi-tli-classical}
  \mathrm{CFI}=1-\frac{\max(0,T_u-df_u)}{\max(0,T_b-df_b)},\qquad
  \mathrm{TLI}=\frac{T_b/df_b-T_u/df_u}{T_b/df_b-1}.
\end{equation}
Under an estimated weight neither $T_u$ nor $T_b$ is $\chi^2$: each is a
mixture $\sum_j\lambda_j\chi^2_1$ whose mean is the generalized degrees of
freedom $\bar Q=\tr(Q\Gamma_x)$, not the nominal $df$. The single new
statistical object is the \emph{joint} law of $(T_u,T_b)$. The headline is that
both statistics are smooth functionals of the \emph{same} $x$, so their joint
covariance is one bilinear form in the \emph{one} shared $\Gamma_x$:
\begin{equation}\label{eq:headline}
  \Cov(T_u,T_b)\;=\;N\,g_u'\,\Gamma_x\,g_b,
\end{equation}
and CFI and TLI are then a ratio delta-method on the two correlated
noncentralities. No new covariance machinery is needed beyond the second
gradient $g_b$ and the second profile Hessian $Q_b$; the metric $\Gamma_x$ is
shared. Three consequences follow: an honest Wald interval for CFI
(Corollary~\ref{cor:cfi}), a TLI interval that is the CFI interval rescaled by a
generalized-df ratio (Corollary~\ref{cor:tli}), and a baseline-dominated
simplification (Remark) under which CFI's variance collapses to the
RMSEA-side variance the program already computes, divided by the squared
baseline noncentrality.

\section{Setup}\label{sec:setup}

Write $x=(u,\gamma)\in\mathbb R^{2m}$ for the extended first-stage vector: the
$m$ polychoric sample moments $u$ (thresholds and correlations, stacked in the
\texttt{OrdinalStats} order) and the DWLS weight diagonal $\gamma=\operatorname{diag}\hat\Gamma$,
$W=\operatorname{diag}(1/\gamma)$. Let
\[
  \sqrt N(\hat x-x_0)\dto\mathcal N(0,\Gamma_x),
\]
with $\Gamma_x$ the joint per-unit NACOV built from the stacked casewise
influence $[\,g_i\mid \mathrm{IF}_i(\gamma)\,]$ exactly as in
\texttt{ordinal\_dwls\_profile\_rmsea}. $\Gamma_x$ is a property of the data, not
of any fitted model: \emph{the user model and the baseline share it.}

For a model $M\in\{u,b\}$ (user, baseline) the minimum-discrepancy fit profiles
its parameters out of
\[
  F_M(x)=d_M(x)'\,W\,d_M(x),\qquad d_M=\sigma_M(\hat\theta_M)-u,
\]
giving statistic $T_M=N F_M$, profile Hessian $Q_M$ (the
$2m\times2m$ extended object of \texttt{weighted\_moment\_profile\_rmsea\_estimated\_weight}),
generalized df $\bar Q_M=\tr(Q_M\Gamma_x)$, and bias-corrected noncentrality
\begin{equation}\label{eq:delta}
  \hat\delta_M \;=\; T_M-\bar Q_M.
\end{equation}
The user model is the fitted CFA/SEM. The baseline is the independence model:
free thresholds, all polychoric correlations fixed to zero (delta scaling fixes
the variances). Its residual is $d_b=(\,0;\,-\operatorname{vech}_{<}\hat R\,)$
(thresholds profile to zero residual under diagonal $W$; the correlation block
is the negative sample correlations), and $T_b$ matches
\texttt{ordinal\_baseline\_chi2}. Because $\sigma_b$ is \emph{linear} in the
baseline thresholds, the baseline observed bread equals its Gauss--Newton bread
and $Q_b$ is exact.

\paragraph{Envelope gradient.} Since $\hat\theta_M$ minimizes $F_M$, the
$\hat\theta_M$-movement term in $dF_M/dx$ vanishes (envelope theorem) and the
gradient is the bare profile score, identical in form for both models:
\begin{equation}\label{eq:grad}
  g_M=\Big(\tfrac{\partial F_M}{\partial u},\ \tfrac{\partial F_M}{\partial\gamma}\Big)
     =\big(-2\,W d_M,\ -d_M^{\,2}/\gamma^{2}\big)\in\mathbb R^{2m},
\end{equation}
the $\gamma$ block coming from $\partial W/\partial\gamma_k=-1/\gamma_k^2$ in
$F_M=\sum_k d_{M,k}^2/\gamma_k$. Equation~\eqref{eq:grad} is exactly the RMSEA
envelope score, evaluated once at the user residual $d_u$ and once at the
baseline residual $d_b$.

\section{Joint law of the two statistics}\label{sec:joint}

\begin{theorem}\label{thm:joint}
Under fixed misspecification of both models and the regularity of the
single-criterion program, with $\hat x=\bar x$ the sample first-stage vector,
\[
  \sqrt N\begin{pmatrix}F_u-F_{0,u}\\[2pt]F_b-F_{0,b}\end{pmatrix}
  \dto
  \mathcal N\!\left(0,\;
  \begin{pmatrix}
    g_u'\Gamma_x g_u & g_u'\Gamma_x g_b\\[2pt]
    g_b'\Gamma_x g_u & g_b'\Gamma_x g_b
  \end{pmatrix}\right).
\]
Equivalently, with $V_{uu}=N g_u'\Gamma_x g_u$, $V_{bb}=N g_b'\Gamma_x g_b$,
$V_{ub}=N g_u'\Gamma_x g_b$, the statistic pair satisfies
$\Var(T_u)=V_{uu}$, $\Var(T_b)=V_{bb}$, $\Cov(T_u,T_b)=V_{ub}$, which is
\eqref{eq:headline}.
\end{theorem}

\begin{proof}
Each $F_M(\hat x)=F_{0,M}+g_M'(\hat x-x_0)+o_p(N^{-1/2})$ by a first-order
expansion, with $g_M$ the envelope gradient \eqref{eq:grad} (the profiled
parameters contribute no first-order term). Stacking the two expansions and
applying the single CLT $\sqrt N(\hat x-x_0)\dto\mathcal N(0,\Gamma_x)$ to the
common $\hat x$ gives the bivariate limit; the off-diagonal is
$g_u'\Gamma_x g_b$ because both rows read the same increment $\hat x-x_0$.
Scaling by $N$ ($T_M=NF_M$) gives the stated variances.
\end{proof}

The mean corrections $\bar Q_M$ are $O_p(1)$ and, as in the single-criterion
program, are treated as fixed in the variance (their sampling fluctuation is
higher order). Hence $\Var(\hat\delta_u)=V_{uu}$, $\Var(\hat\delta_b)=V_{bb}$,
$\Cov(\hat\delta_u,\hat\delta_b)=V_{ub}$.

\section{CFI}\label{sec:cfi}

In the regime where both models misfit, $\hat\delta_u>0$ and $\hat\delta_b>0$,
the truncations in \eqref{eq:cfi-tli-classical} are inactive and, writing the
generalized-df noncentralities for $df$,
\begin{equation}\label{eq:cfi}
  \cfi \;=\; 1-r,\qquad r:=\frac{\hat\delta_u}{\hat\delta_b}.
\end{equation}
The population target is $\mathrm{CFI}_0=1-\delta_{0,u}/\delta_{0,b}
=1-F_{0,u}/F_{0,b}$ (the $N$ cancels in the ratio).

\begin{corollary}\label{cor:cfi}
The ratio gradient is $\nabla r=\hat\delta_b^{-1}(1,-r)$, so by the delta method
\begin{equation}\label{eq:varcfi}
  \widehat\Var(\cfi)=\widehat\Var(r)
  =\frac{1}{\hat\delta_b^{2}}\Big(V_{uu}-2\,r\,V_{ub}+r^{2}V_{bb}\Big),
\end{equation}
and a two-sided $1-\alpha$ interval is
$\cfi\pm z_{1-\alpha/2}\,\widehat\Var(\cfi)^{1/2}$, clipped to $[0,1]$.
\end{corollary}

\section{TLI}\label{sec:tli}

TLI is the same comparison taken per generalized degree of freedom. With
$E[T_M]\approx\bar Q_M+\delta_M$, the misspecification-robust analogue of
$T/df$ is $T_M/\bar Q_M\approx 1+\delta_M/\bar Q_M$, and substituting into the
TLI form of \eqref{eq:cfi-tli-classical} gives
\begin{equation}\label{eq:tli}
  \tli=\frac{T_b/\bar Q_b-T_u/\bar Q_u}{T_b/\bar Q_b-1}
       =1-c\,r,\qquad c:=\frac{\bar Q_b}{\bar Q_u}.
\end{equation}
Thus TLI is CFI's ratio $r$ rescaled by the generalized-df ratio $c$;
$\tli=1-(1-\cfi)\,\bar Q_b/\bar Q_u$. Using $\bar Q$ rather than the nominal
$df$ keeps TLI internally consistent with the mixture mean that defines
$\hat\delta$; when the weight is correctly specified $\bar Q_M\to df_M$ and
\eqref{eq:tli} recovers the lavaan point value.

\begin{corollary}\label{cor:tli}
Treating $c$ as fixed (the $\bar Q_M$ are deterministic mean corrections),
\begin{equation}\label{eq:vartli}
  \tli=1-c\,r\ \Longrightarrow\ \widehat\Var(\tli)=c^{2}\,\widehat\Var(\cfi).
\end{equation}
The TLI interval is the CFI interval scaled by $c=\bar Q_b/\bar Q_u$: same $r$,
same $\widehat\Var(r)$, one extra factor. (TLI is not clipped to $[0,1]$; it can
exceed $1$ for $\hat\delta_u<0$.)
\end{corollary}

\section{The baseline-dominated regime}\label{sec:dominated}

Under any appreciable misspecification the independence model fits grossly, so
$\hat\delta_b\gg\hat\delta_u$, $r\to0$, and \eqref{eq:varcfi} collapses to its
leading term:
\begin{equation}\label{eq:dominated}
  \widehat\Var(\cfi)
  =\frac{V_{uu}}{\hat\delta_b^{2}}\Big(1-2r\,\tfrac{V_{ub}}{V_{uu}}
    +r^2\,\tfrac{V_{bb}}{V_{uu}}\Big)
  \;\approx\;\frac{V_{uu}}{\hat\delta_b^{2}}
  =\frac{\Var(T_u)}{\hat\delta_b^{2}}.
\end{equation}
That is, to leading order CFI's sampling variance is the \emph{user-model}
statistic variance $V_{uu}=N g_u'\Gamma_x g_u$ already returned as
\texttt{grad\_var} by the RMSEA path, divided by the squared baseline
noncentrality $\hat\delta_b^2$; the baseline acts as a nearly constant scale.
The cross-term and baseline-variance contributions are $O(r)$ and $O(r^2)$
corrections. The full bivariate form \eqref{eq:varcfi} is retained in code for
honesty, and it must be: the Monte-Carlo check
(\texttt{tests/checks/ordinal\_cfi\_inference}) finds the ratio
$(V_{uu}/\hat\delta_b^2)/\widehat\Var(\cfi)$ falling from $0.99$ at weak misfit to
$0.35$ at strong misfit, i.e.\ \eqref{eq:dominated} is accurate only near correct
specification, where ``CFI inference is a rescaling of RMSEA inference'' holds;
as misfit grows the cross-term and baseline-variance corrections take over. The
simplification is intuition, not the computation.

\section{Estimated versus fixed weight}\label{sec:weight}

Zeroing the $\gamma$ blocks of $\Gamma_x$ (drop the
$\operatorname{IF}(\gamma)$ rows/columns) gives the fixed-weight comparator: the
$g_M$ $\gamma$-channel in \eqref{eq:grad} is then inert, $V_{MM}$ and $V_{ub}$
revert to their conventional (treat-$W$-as-known) values, and $\bar Q_M$ reverts
to the fixed-weight trace. The estimated-weight law differs from current
practice exactly through this channel. It enters CFI twice: through $V_{uu}$
(the RMSEA-metric piece, where the channel is large because RMSEA's metric
\emph{is} the weight) and through the baseline objects $\hat\delta_b,V_{bb}$.
One might expect the former to dominate and CFI to inherit RMSEA's large weight
sensitivity, but empirically the \emph{net} effect on $\widehat\Var(\cfi)$ is
small: the Monte-Carlo check puts the $\gamma$ share at only $\pm 3\text{--}8\%$
(versus $-12\%$ to $-80\%$ for RMSEA). The ratio structure and the
$\hat\delta_b^2$ scaling attenuate the channel, and the $\hat\delta_b$ shift
under a fixed weight largely cancels it, so CFI sits at the CRMR end of the
spectrum: its estimated- and fixed-weight intervals nearly coincide.

\section{Boundary and regularity}\label{sec:boundary}

\begin{itemize}
\item \textbf{User boundary.} As $\hat\delta_u\downarrow0$ (correct or
near-correct user model) $\cfi\to1$ sits on the parameter-space boundary, the
truncation $\max(0,\cdot)$ activates, and the normal approximation degrades.
The interval is clipped at $1$ and flagged. This is the RMSEA close-fit boundary
made sharper by CFI's hard ceiling.
\item \textbf{Baseline boundary.} If $\hat\delta_b\le0$ (independence model
fits; essentially uncorrelated data) the ratio is undefined: report
$\cfi=1$, $\tli$ and both intervals \texttt{NaN}, with a warning. The delta
method requires $\hat\delta_b$ bounded away from zero, which holds whenever the
baseline is a genuine null.
\item \textbf{Shared metric.} $g_u'\Gamma_x g_b$ uses the \emph{same} $\Gamma_x$
for both gradients; the implementation must assert the two profile fits returned
a common $\Gamma_x$ (they do, being data-only) before forming the cross-term.
\item \textbf{TLI variance is ill-conditioned at strong misfit.} Corollary
\ref{cor:tli} treats $c=\bar Q_b/\bar Q_u$ as fixed, but $\bar Q_u$ is a signed
trace that can be small (and near its sign-change cancellation) when the user
model misfits strongly, giving $c^2$ a heavy right tail. The Monte-Carlo check
finds the TLI interval calibrated at weak/moderate misfit but
\emph{over-covering} at strong misfit (analytic variance $\sim$8$\times$ the
empirical), i.e.\ conservative, not anti-conservative. The TLI \emph{point}
stays unbiased. CFI carries no such division and is the index to trust for an
interval; a stabilized $c$ for TLI is future work.
\item \textbf{Groups.} The pooling is in Section~\ref{sec:multigroup}.
\end{itemize}

\section{Multi-group}\label{sec:multigroup}

For $G$ groups with sizes $n_b$, $N=\sum_b n_b$, the first-stage vectors
$\hat x_b$ are independent across groups, $\sqrt{n_b}(\hat x_b-x_{0,b})\dto
\mathcal N(0,\Gamma_{x,b})$. The fit minimizes the pooled discrepancy
$F=\sum_b (n_b/N)F_b$, and every statistic in this note is a sample-size-weighted
sum of per-group criteria,
\begin{equation}\label{eq:mg-stat}
  T=N\!\sum_b \tfrac{n_b}{N}\,\mathrm{crit}_b=\sum_b n_b\,\mathrm{crit}_b
  \qquad(\mathrm{crit}_b=F_b\ \text{for RMSEA},\ G_b\ \text{for CRMR}).
\end{equation}
Each $\mathrm{crit}_b$ depends on $\hat x_b$ only, with per-group envelope
gradient $g_b$. Because the groups are independent, the joint law collapses to a
\emph{block-diagonal} sum:
\begin{equation}\label{eq:mg-var}
  \Var(T)=\sum_b n_b\,g_b'\,\Gamma_{x,b}\,g_b,\qquad
  \Cov(T^{(1)},T^{(2)})=\sum_b n_b\,g_{b}^{(1)\prime}\,\Gamma_{x,b}\,g_{b}^{(2)},
\end{equation}
the second line being the CFI/TLI cross-term with the user and baseline gradients
(cross-group terms vanish). The profile machinery already returns the block-%
diagonal $\Gamma_x=\operatorname{blkdiag}_b\Gamma_{x,b}$ and the pooled signed
trace $\bar Q=\tr(Q\Gamma_x)$ (the two-metric core scales each block's Jacobian by
$\sqrt{n_b/N}$, giving the correct pooled trace), so $\hat\delta=T-\bar Q$ is
unchanged. The baseline is the per-group independence model: $T_b=\sum_b n_b
F_{b,\mathrm{base}}$, $df_b=\sum_b \mathrm{ncorr}_b$, built as $G$ analytic blocks
through the same primitive with the \emph{shared} per-group $\Gamma_{x,b}$.

In code \eqref{eq:mg-var} is one quadratic form: stack the per-group gradients
into $g_\star$ with block $b$ equal to $\sqrt{n_b/N}\,g_b$, so
$g_\star'\Gamma_x g_\star=\frac1N\Var(T)$ matches the single-group
$g'\Gamma_x g$ exactly (with $n_1=N$ it is the same expression). The
fixed-weight comparator zeros the $\gamma$ coordinates \emph{within each block}
(global indices $[\,\mathrm{off}_b+m_b,\ \mathrm{off}_b+2m_b)$), not a single
trailing $\gamma$ band. CFI, TLI, the ratio delta-method, and the boundary
handling are otherwise identical with the pooled $\hat\delta$'s and
\eqref{eq:mg-var}.

\section{Implementation map}\label{sec:impl}

\begin{enumerate}
\item User profile via \texttt{ordinal\_dwls\_profile\_rmsea} $\Rightarrow$
$Q_u,\Gamma_x,T_u,\bar Q_u$; form $g_u$ from \eqref{eq:grad} with the fitted
residual $d_u$ (as the RMSEA path already does).
\item Baseline profile from the analytic independence block: jacobian
$D_b=[\,I_{nth};0\,]$, residual $d_b$, weight $\gamma$, shared $\Gamma_x$,
Gauss--Newton bread $D_b'WD_b$, $K_b=I$. Route through
\texttt{weighted\_moment\_profile\_rmsea\_estimated\_weight} $\Rightarrow$
$Q_b,T_b,\bar Q_b$; cross-check $T_b$ against \texttt{ordinal\_baseline\_chi2}.
Form $g_b$ from \eqref{eq:grad} with $d_b$.
\item $\hat\delta_u=T_u-\bar Q_u$, $\hat\delta_b=T_b-\bar Q_b$;
$V_{uu}=N g_u'\Gamma_x g_u$, $V_{bb}=N g_b'\Gamma_x g_b$,
$V_{ub}=N g_u'\Gamma_x g_b$.
\item CFI \eqref{eq:cfi}, \eqref{eq:varcfi}; TLI \eqref{eq:tli},
\eqref{eq:vartli}; intervals per Corollaries~\ref{cor:cfi}--\ref{cor:tli} with
the boundary handling of Section~\ref{sec:boundary}.
\item \texttt{estimated\_weight=false} zeros the $\gamma$ blocks of $\Gamma_x$
before steps 3--4 (fixed-weight comparator).
\end{enumerate}

\end{document}
